\section*{Abstract}

This applied research project describes the analysis, design, and implementation of a web-based information system for the comprehensive management of a workshop dedicated to technical support and the repair of Information Technology (IT) equipment. Traditionally, many of these establishments have conducted their administrative operations through manual and fragmented processes, relying on paper records or disconnected office tools. This working approach leads to various operational issues that compromise service quality, such as omissions during device reception, a lack of records regarding physical conditions and included accessories, the preparation of diagnostics and quotes without adequate control over modifications and expiration dates, as well as difficulties in spare parts inventory management. Furthermore, financial vulnerability arises when repairs are initiated without a mandatory upfront deposit. Compounding these challenges is the lack of timely updates for customers regarding the status of their devices, along with the absence of audit logs that allow the identification and review of actions taken by workshop personnel.

In response to this problem, and drawing upon the quality management principles established in the ISO 9001:2015 standard as well as scientific literature regarding the automation of maintenance and repair processes for electronic equipment, the purpose of this project is to centralize, automate, and preserve a traceable record of each stage of technical service. The general objective is to develop a web system for controlling and administering an IT equipment repair workshop, utilizing the Laravel framework (PHP 8+) and the Model-View-Controller (MVC) architectural pattern. This solution aims to optimize work order tracking, inventory administration, and customer communication. To achieve this, operational requirements and workshop business rules were analyzed, the software architecture and relational database model were designed, and an interactive dashboard was implemented to monitor repair statuses in real time.

The system incorporates a security framework based on Role-Based Access Control (RBAC), structured into five distinct user profiles: Administrator, Reception, Technician, Warehouse, and Customer. Additionally, it integrates functional modules for the formal reception of computers, printers, and peripherals through the automatic generation of unique tracking numbers (folios), accessory logging, and photographic evidence storage; technical diagnostics and quote formulation with version control and validity periods; and customer approval or rejection workflows prior to starting any physical repair. The platform also encompasses inventory management with dedicated spare parts reservation linked to authorized service orders, alongside the registration of internal tool loans. In the financial domain, it manages payments by enforcing a mandatory 50\,\% deposit, records final settlements, and applies loyalty discounts. Furthermore, it embeds quality control procedures through technical testing prior to delivery, technical warranty issuance linked directly to original service orders, and comprehensive tracking of user activities through an immutable audit trail.

With the implementation of this platform, the workshop aims to enhance productivity and profitability, prevent financial and material losses caused by misplaced devices or components, and reduce repair turnaround times. Likewise, it seeks to strengthen customer trust by providing an accessible web platform that enables users to transparently monitor the progress of their services and review relevant technical documentation from any internet-connected device.

\vspace{0.5cm}
\noindent\textbf{Keywords:} Service order management, IT equipment maintenance, Laravel, Operational traceability, Inventory control, Technical warranties, Role-Based Access Control (RBAC).
